\chapter{Geradores de analisadores léxicos}

\section{Introdução}

Neste capítulo, veremos como utilizar um gerador para construir um
analisador léxico a partir de uma especificação. Em Haskell, o
gerador de analisadores mais popular é o Alex. Este capítulo
apresenta uma introdução ao uso do Alex em um exemplo prático: a
construção do analisador léxico para uma linguagem de expressões.

O código discutido neste capítulo esta disponível em
\begin{center}
  \verb|Exp/Frontend/Lexer/Alex/ExpLexer.x|
\end{center}

\section{Construindo um analisador léxico utilizando o Alex}

Nesta seção vamos apresentar a especificação de um analisador léxico
utilizando a linguagem do Alex. Primeiro, vamos definir um tipo para
representação de tokens e, em seguida, descreveremos a especificação
do analisador e como este pode ser criado utilizando o Alex.

\subsection{Definindo os tokens}

Antes de criar o lexer, precisamos definir o que é um \textbf{token}
na nossa linguagem de expressões. Os tipos a seguir apresentam a
definição de token que utilizaremos.

\begin{minted}{haskell}
data Token = Token {
      pos    :: (Int, Int) -- (Linha, Coluna)
    , lexeme :: Lexeme
    } deriving (Eq, Show)

data Lexeme
  = TNumber Int | TLParen | TRParen | TPlus | TTimes | TEOF
  deriving (Eq, Show)
\end{minted}

O tipo \haskell{Lexeme} representa as diferentes
categorias de tokens presentes na linguagem: números, abre e fecha
parêntesis, adição, multiplicação e um valor especial para fim da
entrada (\haskell{EOF}). Um token é formado por um
\haskell{Lexeme} e um par de inteiros que representam
a linha e coluna inicial daquele token na entrada.

\subsection{Sintaxe de expressões regulares no Alex}

Alex usa expressões regulares para definir padrões. A seguir, vamos
apresentar alguns exemplos de expressões que são utilizadas na
especificação léxica de linguagens.

A sintaxe \verb|[0-9]| denota que qualquer
caractere numérico entre 0 e 9 (inclusive) será aceito. Da mesma
forma, \verb|[a-zA-Z]| denota uma expressão para
qualquer caractere minúsculo ou maiúsculo. Caracteres também podem
ser excluídos; por exemplo, \verb|[\^\\"]| é
entendido como ``qualquer coisa, exceto aspas duplas''.

O operador Kleene denota que uma expressão pode ser casada zero ou
mais vezes com a entrada. Da mesma forma, o fecho positivo,
\verb|+| denota que uma expressão pode ser casada
uma ou mais vezes. Por exemplo, \verb|[0-9]+|
aceita a palavra 1234, pois denota uma sequência formada por um ou
mais dígitos. Outro operador comum é \verb|?| que
significa que uma expressão pode ser casada opcionalmente, isto é,
zero ou uma vez.

Um ponto, \verb|.|, casa com qualquer caractere
diferente de quebra de linha. Uma string entre aspas duplas, como
``>='', corresponderá exatamente à string dentro das
aspas. Um caractere de escape, como \verb|\n|, corresponderá exatamente a
esse caractere.

Você pode agrupar expressões regulares concatenando-as.
Ao concatenar expressões
regulares, como em \verb|[a-z][A-Za-z0-9]*|,
significa que o Alex deve corresponder a cada conjunto de strings em
sequência. Um exemplo de lexema que pode ser aceita por essa
expressão regular é \verb|myVariable1|.

O operador de barra vertical, \verb|||, indica que o Alex pode
alternar entre opções, como

\begin{verbatim}
0(x[0-9a-fA-F]+ | o[0-7]+)
\end{verbatim}

que é uma expressão regular que pode corresponder a números
hexadecimais ou octais.

\subsection{Estrutura de especificações Alex}

A estrutura de um arquivo de especificação Alex (com extensão .x) é
dividida em quatro seções principais. O arquivo inicia com um bloco
de código Haskell, delimitado por chaves $\{ \ldots{} \}$, onde são
inseridos o nome do módulo e as importações necessárias. No exemplo,
\verb|ExpLexer.x|, temos o seguinte bloco inicial:

\begin{minted}{haskell}
{
{-# OPTIONS_GHC -Wno-name-shadowing #-}
module Exp.Frontend.Lexer.Alex.ExpLexer where

import Control.Monad
import Exp.Frontend.Lexer.Token
}
\end{minted}

Logo abaixo, encontram-se as definições de macros e diretivas de
configuração, como o \verb|%wrapper|, que instrui o Alex sobre qual o
``template'' de analisador léxico deve ser gerado.

\begin{verbatim}
%wrapper "monadUserState"

$digit = 0-9
@number     = $digit+
\end{verbatim}

O Alex suporta dois tipos de macros para facilitar a escrita de
expressões regulares: os conjuntos de caracteres (como
\verb|$digit|), que representam classes de
caracteres, e as macros de expressões regulares (como
\verb|@number|), que funcionam como atalhos para
padrões mais complexos. Quanto à infraestrutura de código gerada, o
Alex utiliza wrappers que definem a interface do lexer; os mais
comuns incluem o \textbf{basic}, que processa uma String de forma
simples, o \textbf{posn}, que adiciona automaticamente informações de
linha e coluna, e o \textbf{monadUserState}, que permite a execução
dentro de uma mônada para manipulação de estados complexos. A
principal diferença entre eles reside no controle: enquanto os
wrappers simples são puramente funcionais e diretos, o wrapper
\textbf{monadUserState} (utilizado no exemplo) oferece maior
flexibilidade ao permitir que o desenvolvedor armazene dados
personalizados, como o nível de aninhamento de comentários,
diretamente no estado do analisador.

A especificação léxica é inicializada pelo símbolo
\verb|tokens :-| e contém as regras de tradução
que associam expressões regulares a ações específicas em Haskell, que
podem ser utilizadas para criação de tokens da linguagem. A seguir,
apresentamos a especificação para a linguagem de expressões.

\begin{verbatim}
tokens :-
      -- whitespace and line comments
      <0> $white+       ;
      <0> "//" .*       ;
      -- other tokens
      <0> @number       {mkNumber}
      <0> "("           {simpleToken TLParen}
      <0> ")"           {simpleToken TRParen}
      <0> "+"           {simpleToken TPlus}
      <0> "*"           {simpleToken TTimes}
      -- multi-line comment
      <0> "\*"              { nestComment `andBegin` state_comment }
      <0> "*/"              {\ _ _ -> alexError "Error! Unexpected close comment!" }
      <state_comment> "\*"  { nestComment }
      <state_comment> "*/"  { unnestComment }
      <state_comment> .     ;
      <state_comment> \n    ;
\end{verbatim}

O código anterior define as regras de reconhecimento para cada
elemento da linguagem, utilizando \textbf{start codes} para alternar
entre o processamento normal e o de comentários. No estado inicial
\verb|<0>|, o analisador ignora espaços em branco
(\verb|$white+|) e comentários de linha (começam
com \verb|//|). Para tokens significativos como
números (\verb|@number|), parênteses e operadores,
o Alex executa ações Haskell como \verb|mkNumber|
ou \verb|simpleToken|, que geram os tokens correspondentes. Caso o
analisador encontre um token para fechar um comentário,
\verb|*/| enquanto ainda está no estado inicial,
ele dispara um erro imediato, pois não há um bloco aberto para ser fechado.

A transição para o estado de comentários multilinha ocorre quando o
padrão \verb|\\*| é identificado no estado
\verb|<0>|, acionando a função
\verb|nestComment| e mudando o estado para
\verb|state_comment| através da função
\verb|andBegin|. Uma vez dentro do estado
\verb|state_comment|, o léxico ignora caracteres
genéricos (usando \verb|.|) e quebras de linha
(\verb|\\n|), mas permanece atento a novos
símbolos de abertura ou fechamento de blocos de comentário. O uso de
\verb|nestComment| e
\verb|unnestComment| dentro desse estado permite
que o analisador incremente ou decremente um contador de nível de
aninhamento. Somente quando o nível de aninhamento retorna a zero é
que a função unnestComment retorna para o estado
\verb|<0>|, permitindo que a análise continue normalmente.

Em especificações Alex, quando utilizamos o wrapper
\verb|monadUserState|, ganhamos a capacidade de
carregar um estado personalizado durante toda a análise, definido
pelo tipo \verb|AlexUserState|. No exemplo, esse
tipo contém apenas o campo \verb|nestLevel| para
rastrear o nível de blocos de comentários aninhados. Para definir o
estado inicial, devemos codificar a função
\verb|alexInitUserState| que inicializa esse
contador de nível de comentários em zero.

A manipulação desse estado ocorre através de três funções
fundamentais: \verb|get|,
\verb|put| e \verb|modify|. A
função \verb|get| extrai o estado atual de dentro
da mônada Alex, enquanto \verb|put| sobrescreve o
estado existente por um novo. Já a função
\verb|modify| aplica uma transformação direta ao
estado. Por exemplo, na função \verb|nestComment|,
\verb|modify| é usada para incrementar o nível de
comentário de forma concisa.

\section{Conclusão}

Este capítulo apresentou como usar o gerador de analisadores léxicos
para Haskell, Alex. Apresentamos a definição de tokens, como tipos de
dados Haskell. Utilizando a linguagem de especificação do Alex,
definimos quais padrões (definidos como expressões regulares) e quais
tokens são associados a eles. Apresentamos os conceitos de wrappers e
como estados podem ser utilizados para modificar o comportamento do
analisador léxico. Mostramos como esse mecanismo é utilizado para
descartar comentários em bloco. Com isso, temos a base necessária
para a próxima fase do compiladores: a análise sintática.
